\documentclass[10pt]{article}
\usepackage[margin=0.78in]{geometry}
\usepackage[T1]{fontenc}
\usepackage{lmodern}
\usepackage{amsmath,amssymb,amsthm}
\usepackage{booktabs}
\usepackage{microtype}
\usepackage[hidelinks,bookmarks=false]{hyperref}
\newtheorem{theorem}{Theorem}
\newtheorem{proposition}{Proposition}
\newcommand{\Z}{\mathbb Z}
\newcommand{\Prb}{\mathbb P}
\newcommand{\PhiC}{\Phi}
\title{Möbius Reliability on the Twenty-Subgroup Lattice of a Named Abelian Core}
\author{Anonymous}
\date{}

\begin{document}
\maketitle

\begin{abstract}
We compute the exact generating-closure reliability law for sixteen frozen
named coordinates in $Q\cong\Z/9\oplus\Z/3\oplus\Z/2$.  A label is deleted
independently with probability $q$, and the retained labels generate a
subgroup of the actual twenty-node subgroup lattice.  For each subgroup $H$,
we count the labels whose cyclic closure lies in $H$, form the cumulative
probability $P_{\leq H}(q)=q^{16-n_H}$, and apply the integer Möbius function
of that lattice.  An independent enumeration of all $2^{16}=65536$ supports
agrees coefficient-by-coefficient for every subgroup.  The top row is
\(1-q-q^4+q^5-q^7-q^8+5q^9-3q^{10}\), which factors as
\((1-q)(1-q^4-q^7-2q^8+3q^9)\) and reproduces the preceding C73
reliability result.  The calculation is finite and source-bound; it makes no
arithmetic, local, Euler-factor, root-number, Burnside-ring, or Hilbert--Polya
claim.
\end{abstract}

\section{Named finite probability object}
Let
\[
 Q=\Z/9\oplus\Z/3\oplus\Z/2,
 \qquad L=\{S_1,\ldots,S_{16}\}.
\]
The sixteen coordinates and the ordered list of all twenty subgroups are
inherited byte-for-byte from the C76 closure-orbit evidence.  If
$A\subseteq L$ is the retained support, write
\[
 \PhiC(A)=\langle x_i:S_i\in A\rangle\leq Q,
 \qquad \PhiC(\varnothing)=\{0\}.
\]
Every label is retained with probability $1-q$ and deleted with probability
$q$.  Thus the exact probability of a closure event is a polynomial in the
single formal variable $q$; no floating-point evaluation is used.

For each label let $C_i=\langle x_i\rangle$.  Define
\[
 n_H=\#\{i:C_i\leq H\}.
\]
The labels outside this set must be deleted for the generated subgroup to be
contained in $H$, while labels in the set are unrestricted.

\begin{theorem}[Cumulative law]
For every one of the twenty source-ordered subgroups,
\[
 P_{\leq H}(q):=\Prb[\PhiC(A)\leq H]=q^{16-n_H}.
\]
\end{theorem}

The C73--C76 source binding is
\[
\begin{split}
\operatorname{sha256}(\text{C73 evidence})&=\texttt{e91c8e6dcf1de5362b1a052ada83eb758b2c2d75520c1e8bdbd37ab055c725e5},\\
\operatorname{sha256}(\text{C73 manifest})&=\texttt{a50b5707d36f8b94b463e6c5fc4b5b7f6d6df7eb5e87d70bfc82d2b1a653cd8d},\\
\operatorname{sha256}(\text{C75 evidence})&=\texttt{8beee17a227153e066907549df70c14a087b7de4141c3092d7cebd4a91541d98},\\
\operatorname{sha256}(\text{C75 manifest})&=\texttt{7ede3e35c3101d17c683d2da440037d5bd4e002266530b52b3d1cb36ed4c8fcb},\\
\operatorname{sha256}(\text{C76 evidence})&=\texttt{42e7783b2652666b84ea7f82b65d2421d98064ee5d5011ab94033aa18c051a94},\\
\operatorname{sha256}(\text{C76 manifest})&=\texttt{55725664005113ae993b54197ff4fbd97bde347ce49aa69ea0c228372ba289d5}.
\end{split}
\]

\section{Möbius inversion}
Let $K\leq H$ denote containment in the actual subgroup poset, and let
$\mu(K,H)$ be its integer Möbius function.  The cumulative events satisfy
$P_{\leq H}=\sum_{K\leq H}P_{=K}$, so inversion gives the exact identity
\[
 P_{=H}(q)=\sum_{K\leq H}\mu(K,H)q^{16-n_K}.
 \tag{1}
\]

The source-ordered subgroup orders and containment counts are
{\scriptsize
\[
\begin{array}{c|rrrrrrrrrrrrrrrrrrrr}
 i&0&1&2&3&4&5&6&7&8&9&10&11&12&13&14&15&16&17&18&19\\\hline
 |H_i|&1&2&3&3&3&3&6&6&6&6&9&9&9&9&18&18&18&18&27&54\\
 n_{H_i}&5&6&7&6&7&6&8&7&8&7&11&7&7&8&12&8&8&9&15&16
\end{array}
\]
}\normalsize

For a direct check, let $c_{H,k}$ be the number of supports of size $k$ with
closure exactly $H$.  Exhaustive enumeration yields
\[
 P^{\rm dir}_{=H}(q)=\sum_{k=0}^{16}c_{H,k}(1-q)^kq^{16-k}.
 \tag{2}
\]

\begin{proposition}[Exact inversion check]
For all twenty subgroups, the integer coefficient vectors obtained from (1)
and (2) are identical.  The direct support totals by index are
{\scriptsize
\[
\begin{array}{c|rrrrrrrrrrrrrrrrrrrr}
 i&0&1&2&3&4&5&6&7&8&9&10&11&12&13&14&15&16&17&18&19\\\hline
 \sum_kc_{H_i,k}&32&32&96&32&96&32&96&32&96&32&1760&64&64&192&1760&64&64&192&30400&30400.
\end{array}
\]
}\normalsize
These totals sum to $65536$, and $\sum_H P_{=H}(q)=1$ identically.
\end{proposition}

\section{The complete polynomial atlas}
For compactness, each row below lists the nonzero terms as
``exponent: coefficient'' pairs in $q$.  This is the canonical integer
coefficient representation of (1), not a numerical approximation.

\small
\begin{center}
\begin{tabular}{c|l}
index $i$ & nonzero coefficients of $P_{=H_i}(q)$\\\hline
0 & $11:1$\\
1 & $10:1,\ 11:-1$\\
2 & $9:1,\ 11:-1$\\
3 & $10:1,\ 11:-1$\\
4 & $9:1,\ 11:-1$\\
5 & $10:1,\ 11:-1$\\
6 & $8:1,\ 9:-1,\ 10:-1,\ 11:1$\\
7 & $9:1,\ 10:-2,\ 11:1$\\
8 & $8:1,\ 9:-1,\ 10:-1,\ 11:1$\\
9 & $9:1,\ 10:-2,\ 11:1$\\
10 & $5:1,\ 9:-2,\ 10:-2,\ 11:3$\\
11 & $9:1,\ 10:-1$\\
12 & $9:1,\ 10:-1$\\
13 & $8:1,\ 10:-1$\\
14 & $4:1,\ 5:-1,\ 8:-2,\ 10:5,\ 11:-3$\\
15 & $8:1,\ 9:-2,\ 10:1$\\
16 & $8:1,\ 9:-2,\ 10:1$\\
17 & $7:1,\ 8:-1,\ 9:-1,\ 10:1$\\
18 & $1:1,\ 5:-1,\ 8:-1,\ 9:-2,\ 10:3$\\
19 & $0:1,\ 1:-1,\ 4:-1,\ 5:1,\ 7:-1,\ 8:-1,\ 9:5,\ 10:-3$
\end{tabular}
\end{center}
\normalsize

The final row corresponds to $H_{19}=Q$:
\[
\begin{split}
 P_{=Q}(q)&=1-q-q^4+q^5-q^7-q^8+5q^9-3q^{10}\\
 &= (1-q)(1-q^4-q^7-2q^8+3q^9).
\end{split}
\tag{3}
\]
Equation (3) is exactly the homogeneous reliability polynomial reported by
C73.  The agreement is a cross-generation consistency check: C77 derives it
from the full subgroup lattice, while C73 obtained it from the generation
criterion.

\section{Reproducibility and scope}
The producer binds all C73, C75, and C76 bytes before reading source data.  An independent
checker reconstructs cyclic closures, containment, the Möbius recursion, and
all 65536 direct supports.  Separate symbolic cross-check, clean replay, and
hostile semantic mutations are required release gates.  The effective C76
label action and its six-element ambient-kernel boundary are retained; no
non-faithful ambient lift is substituted.

This theorem note concerns only a finite named presentation and its deletion
probability polynomials.  It does not claim a full table of marks or Burnside
ring, an arithmetic or local object, Euler factors, root numbers, automorphy,
or a Hilbert--Polya operator.  The scope firewall is
\texttt{NO\_BAD\_EULER\_OR\_ROOT\_NUMBER}.

\end{document}
